% © 2026 Ing. David Jaroš — CC BY-NC-ND 4.0
%
% This work is licensed under a Creative Commons Attribution-NonCommercial-NoDerivatives
% 4.0 International License (CC BY-NC-ND 4.0).
%
% License History: Earlier drafts (up to v0.3) were released under CC BY 4.0.
% From v0.4 onward, all material is released under CC BY-NC-ND 4.0 to protect
% the integrity of the theoretical work during ongoing academic development.
%
% See LICENSE.md for full license text.

% UBT-AI-PROVENANCE-BEGIN
% schema: ubt-ai-provenance/v1
% tier: C_working
% ai_assistance: disclosed
% human_review: risk-based
% editorial_responsibility: Ing. David Jaroš
% policy: ../../AI_PROVENANCE.md
% notice: Working material; exhaustive human review is not claimed.
% UBT-AI-PROVENANCE-END

\ifdefined\INCLUDEMODE
\else
  \documentclass[12pt,a4paper]{article}
  \usepackage[a4paper, margin=2.5cm]{geometry}
  \usepackage{amsmath, amssymb, amsthm}
  \usepackage{mathtools}
  \usepackage{hyperref}
  \usepackage{xcolor}
  \usepackage{mdframed}
  \usepackage{booktabs}
  \usepackage{longtable}
  \usepackage{tcolorbox}

  \newtheorem{theorem}{Theorem}[section]
  \newtheorem{lemma}[theorem]{Lemma}
  \newtheorem{proposition}[theorem]{Proposition}
  \newtheorem{corollary}[theorem]{Corollary}
  \theoremstyle{definition}
  \newtheorem{definition}[theorem]{Definition}
  \theoremstyle{remark}
  \newtheorem{remark}[theorem]{Remark}

  \newmdenv[backgroundcolor=yellow!20, linecolor=orange, linewidth=1.5pt]{gapbox}
  \newmdenv[backgroundcolor=green!15, linecolor=green!60!black, linewidth=1.5pt]{resultbox}
  \newmdenv[backgroundcolor=red!8, linecolor=red!60, linewidth=1.5pt]{warnbox}
  \newmdenv[backgroundcolor=blue!8, linecolor=blue!50, linewidth=1.5pt]{insightbox}

  \title{\textbf{$N_{\mathrm{eff}} = 12$: Five Independent First-Principles Derivations}\\[6pt]
         \large Track E2 — T3\_ALPHA Research}
  \author{Ing.\ David Jaro\v{s}}
  \date{2026-04-28}

  \begin{document}
  \maketitle
  \tableofcontents
  \bigskip
\fi

%──────────────────────────────────────────────────────────────────────────────
\section{Purpose and Context}
\label{sec:neff12:purpose}
%──────────────────────────────────────────────────────────────────────────────

\begin{tcolorbox}[colback=blue!4!white,colframe=blue!50!black,title=Track E2 Objective]
Provide five independent first-principles derivations of $N_{\mathrm{eff}} = 12$
from the UBT axiom $\mathcal{B} = \mathbb{C}\otimes\mathbb{H}$, without using
$\alpha$ or $m_e$ as input.

\medskip
\begin{tabular}{ll}
Companion track & E1: Origins of exponent $3/2$ \\
Companion files & \texttt{exponent\_3\_over\_2\_candidates.tex},
                  \texttt{ALPHA\_STRUCTURAL\_ORIGINS.md} \\
\end{tabular}
\end{tcolorbox}

The effective number of charged modes $N_{\mathrm{eff}} = 12$ enters the derivation
of $\alpha^{-1}_{\mathrm{bare}} = 137$ in two places:

\begin{enumerate}
  \item \textbf{One-loop baseline}:
        $B_0 = \dfrac{2\pi N_{\mathrm{eff}}}{3} = 8\pi$ \quad (proved [L1])

  \item \textbf{Full one-loop coefficient}:
        $B_{\mathrm{base}} = N_{\mathrm{eff}}^{3/2} \approx 41.57$ \quad (conditional [MC])
\end{enumerate}

The value $N_{\mathrm{eff}} = 12$ was previously derived by a single algebraic route
(dimensional decomposition).  This document provides \emph{four additional independent
routes}, all yielding $N_{\mathrm{eff}} = 12$.  The over-determination reinforces the
non-circular nature of the result.

%──────────────────────────────────────────────────────────────────────────────
\section{Route R1: Dimensional Decomposition of $\mathcal{B}$ (Existing)}
\label{sec:neff12:R1}
%──────────────────────────────────────────────────────────────────────────────

This is the existing derivation from~\texttt{canonical/n\_eff/step1\_mode\_decomposition.tex}.
It is reproduced here as the reference derivation.

\begin{theorem}[$N_{\mathrm{eff}} = 12$ from dimensional decomposition]
\label{thm:neff12:R1}
The effective number of distinct charged modes of $\Theta \in \mathcal{B} = \mathbb{C}\otimes\mathbb{H}$
contributing to the one-loop vacuum polarisation of the $U(1)_{\mathrm{EM}}$ gauge
field on $S^1_\psi$ is:
\begin{equation}
  N_{\mathrm{eff}}
  = \underbrace{N_{\mathrm{phases}}}_{=\,\dim_\mathbb{R}(\mathrm{Im}\,\mathbb{H})\,=\,3}
    \times
    \underbrace{N_{\mathrm{helicity}}}_{=\,2}
    \times
    \underbrace{N_{\mathrm{charge}}}_{=\,2}
  = 3 \times 2 \times 2
  = 12.
  \label{eq:neff12:R1}
\end{equation}
\end{theorem}

\begin{proof}
\begin{enumerate}
  \item \textbf{Phase factor}: The three imaginary quaternion generators
    $\mathbf{i}, \mathbf{j}, \mathbf{k} \in \mathrm{Im}\,\mathbb{H}$ are
    algebraically independent.  Each generates an independent $U(1)$ phase
    action on $\Theta$.  The number of independent phase directions is
    $\dim_\mathbb{R}(\mathrm{Im}\,\mathbb{H}) = 3$.

  \item \textbf{Helicity factor}: The complexification $\mathbb{C}\otimes$ provides
    a complex structure on $\mathbb{H}$, splitting each quaternion field component
    into left and right Weyl helicity sectors.  Thus $N_{\mathrm{helicity}} = 2$.

  \item \textbf{Charge factor}: The complex structure of $\mathcal{B}$ provides
    a charge-conjugation involution $\Theta \mapsto \Theta^*$, doubling the
    spectrum into particle and antiparticle modes.  Thus $N_{\mathrm{charge}} = 2$.

  \item \textbf{Independence}: The three factors count distinct quantum numbers
    ($U(1)$ phase direction, helicity, charge sign) that commute.  They multiply.
\end{enumerate}
\end{proof}

\begin{resultbox}
Status: CLEAN [L0].  Zero free parameters.  No reference to $\alpha$ or $m_e$.\\
Source: \texttt{canonical/n\_eff/step1\_mode\_decomposition.tex}
\end{resultbox}

%──────────────────────────────────────────────────────────────────────────────
\section{Route R2: Standard Model Generator Count}
\label{sec:neff12:R2}
%──────────────────────────────────────────────────────────────────────────────

\subsection{SM Gauge Group Embedding}
\label{ssec:neff12:R2:embedding}

The Standard Model gauge group $G_{\mathrm{SM}} = SU(3)_c \times SU(2)_L \times U(1)_Y$
is derived from $\mathcal{B} = \mathbb{C}\otimes\mathbb{H}$ in:
\begin{center}
\texttt{canonical/interactions/sm\_gauge.tex}
\end{center}

The embedding produces a one-to-one correspondence between generators of $G_{\mathrm{SM}}$
and independent gauge-charged modes of the $\Theta$-field that contribute to
$U(1)_{\mathrm{EM}}$ vacuum polarisation.

\subsection{Generator Counting}
\label{ssec:neff12:R2:count}

\begin{theorem}[$N_{\mathrm{eff}} = 12$ from SM generators]
\label{thm:neff12:R2}
The total number of generators of $G_{\mathrm{SM}}$ embedded in $\mathcal{B}$ is:
\begin{equation}
  \dim G_{\mathrm{SM}}
  = \underbrace{8}_{SU(3)_c}
  + \underbrace{3}_{SU(2)_L}
  + \underbrace{1}_{U(1)_Y}
  = 12.
  \label{eq:neff12:R2}
\end{equation}
Each generator corresponds to one independent gauge-charged mode; hence
$N_{\mathrm{eff}} = \dim G_{\mathrm{SM}} = 12$.
\end{theorem}

\begin{proof}
$SU(3)$ has $3^2 - 1 = 8$ generators (Gell-Mann matrices $\lambda_a$, $a=1,\ldots,8$).
$SU(2)$ has $2^2 - 1 = 3$ generators (Pauli matrices $\tau_a$, $a=1,2,3$).
$U(1)$ has $1$ generator (hypercharge $Y$).
Total: $8 + 3 + 1 = 12$.

Each generator $T^a$ of $G_{\mathrm{SM}}$ couples to the $\Theta$-field via
the covariant derivative $D_\mu = \partial_\mu + igA^a_\mu T^a$.  The virtual
loop correction to the $U(1)_{\mathrm{EM}}$ propagator receives one independent
contribution per generator (in the adjoint representation).
\end{proof}

\begin{remark}[Relationship to Route R1]
The generator count $8+3+1=12$ agrees with $3\times2\times2=12$ because the SM
gauge group structure in the biquaternion algebra satisfies:
$\dim SU(3) + \dim SU(2) + \dim U(1) = N_{\mathrm{phases}} \times (N_{\mathrm{helicity}} \times N_{\mathrm{charge}}) = 3 \times 4 = 12$.
This is a consistency check.
\end{remark}

\begin{resultbox}
Status: CLEAN [L0].  The SM generator count is derived from $\mathcal{B}$ in
\texttt{sm\_gauge.tex}; the route is independent of Route~R1.\\
Independence: $\dim G_{\mathrm{SM}} = 12$ is a pure group-theoretic fact;
no reference to $\alpha$ or $m_e$.
\end{resultbox}

%──────────────────────────────────────────────────────────────────────────────
\section{Route R3: Three-Qubit Internal Sector Decomposition}
\label{sec:neff12:R3}
%──────────────────────────────────────────────────────────────────────────────

\subsection{Qubit Encoding of the SM Gauge Charges}
\label{ssec:neff12:R3:encoding}

The colour, isospin, and hypercharge quantum numbers of the SM can be encoded
in a 3-qubit register (see~\texttt{canonical/interactions/su3\_qubit\_encoding.tex}):

\begin{center}
\begin{tabular}{lll}
\toprule
Sector & Quantum number & Qubit dimension \\
\midrule
Colour ($SU(3)_c$) & $r, g, b$ & 3 (one-hot: $|r\rangle, |g\rangle, |b\rangle$) \\
Isospin ($SU(2)_L$) & up, down & 2 ($|\uparrow\rangle, |\downarrow\rangle$) \\
Hypercharge ($U(1)_Y$) & $Y = \pm 1/2$ & 1 ($|+\rangle, |-\rangle$) \\
\bottomrule
\end{tabular}
\end{center}

The total number of independent internal charged modes before charge-conjugation is:
\begin{equation}
  N_{\mathrm{internal}} = 3 + 2 + 1 = 6.
  \label{eq:neff12:R3:internal}
\end{equation}

\subsection{Charge-Conjugation Doubling}
\label{ssec:neff12:R3:doubling}

\begin{theorem}[$N_{\mathrm{eff}} = 12$ from 3-qubit sectors]
\label{thm:neff12:R3}
Including charge-conjugation (particle/antiparticle doubling from the complex
structure $\mathbb{C}\otimes$ in $\mathcal{B}$):
\begin{equation}
  N_{\mathrm{eff}}
  = N_{\mathrm{internal}} \times N_{\mathrm{charge}}
  = (3 + 2 + 1) \times 2
  = 6 \times 2
  = 12.
  \label{eq:neff12:R3}
\end{equation}
\end{theorem}

\begin{proof}
Each internal charged mode (colour, isospin, or hypercharge state) comes in a
particle/antiparticle pair due to the complex structure of $\mathcal{B}$:
if $\Theta$ is a solution, so is $\Theta^\dagger$ with opposite charges.
The 6 internal modes therefore split into 12 distinct modes.
\end{proof}

\begin{remark}[Comparison with Route R2]
In Route R2, the generators are: 8 colour gluons + 3 $W$-bosons + 1 $B$-boson = 12.
The qubit sectors in Route R3 give: 3 quarks + 2 isospin + 1 hypercharge = 6 modes,
times 2 = 12.  The factor of 2 difference between 6 and 12 corresponds exactly
to the charge-conjugation doubling, which is the ``$N_{\mathrm{charge}} = 2$''
factor in Route~R1.
\end{remark}

\begin{resultbox}
Status: CLEAN [L0].  Derived from $\mathcal{B}$ qubit structure.\\
Source: \texttt{canonical/interactions/su3\_qubit\_encoding.tex},
\texttt{canonical/bridges/su3\_gauge\_qubit\_equivalence.tex}\\
Independence: No reference to $\alpha$ or $m_e$.
\end{resultbox}

%──────────────────────────────────────────────────────────────────────────────
\section{Route R4: Off-Diagonal Spinor Component Counting}
\label{sec:neff12:R4}
%──────────────────────────────────────────────────────────────────────────────

\subsection{$\Theta$ as a $2\times 2$ Complex Matrix}
\label{ssec:neff12:R4:matrix}

Since $\mathcal{B} = \mathbb{C}\otimes\mathbb{H} \cong M_2(\mathbb{C})$,
the field $\Theta \in \mathcal{B}$ is representable as a $2\times 2$ complex matrix:
\begin{equation}
  \Theta = \begin{pmatrix} \theta_{11} & \theta_{12} \\ \theta_{21} & \theta_{22} \end{pmatrix},
  \qquad \theta_{ij} \in \mathbb{C}.
  \label{eq:neff12:R4:matrix}
\end{equation}

Under the $U(1)_{\mathrm{EM}}$ gauge action, the four entries transform differently:

\begin{center}
\begin{tabular}{lll}
\toprule
Entry & Gauge character & Physical interpretation \\
\midrule
$\theta_{11}$ & Neutral (singlet) & Right-handed neutral component \\
$\theta_{22}$ & Neutral (singlet) & Left-handed neutral component \\
$\theta_{12}$ & Charged $+q$ & Positive-charge component \\
$\theta_{21}$ & Charged $-q$ & Negative-charge component \\
\bottomrule
\end{tabular}
\end{center}

The two off-diagonal entries $\theta_{12}$ and $\theta_{21}$ are charged.

\subsection{Mode Counting with Quaternion Phases and Helicity}
\label{ssec:neff12:R4:counting}

\begin{theorem}[$N_{\mathrm{eff}} = 12$ from off-diagonal spinor components]
\label{thm:neff12:R4}
The number of independent charged modes contributing to vacuum polarisation is:
\begin{equation}
  N_{\mathrm{eff}}
  = \underbrace{2}_{\text{off-diag entries}}
    \times
    \underbrace{3}_{N_{\mathrm{phases}}\text{ in Im}\,\mathbb{H}}
    \times
    \underbrace{2}_{N_{\mathrm{helicity}}}
  = 2 \times 3 \times 2
  = 12.
  \label{eq:neff12:R4}
\end{equation}
\end{theorem}

\begin{proof}
Each off-diagonal matrix entry $\theta_{12}$ and $\theta_{21}$ carries:
\begin{itemize}
  \item Three independent phase directions from Im$\,\mathbb{H}$
    ($\phi_\mathbf{i}, \phi_\mathbf{j}, \phi_\mathbf{k}$), as each
    $\mathbb{H}$-valued entry is a full quaternion.
  \item Two Weyl helicity components (left/right) from the $\mathbb{C}$ tensor
    factor.
\end{itemize}
The two off-diagonal entries differ by charge sign and are independent.
Total independent charged modes: $2 \times 3 \times 2 = 12$.
\end{proof}

\begin{remark}
The neutral diagonal entries $\theta_{11}, \theta_{22}$ do not contribute to
$U(1)_{\mathrm{EM}}$ vacuum polarisation at one loop.  Excluding them gives the
same count as including them but subtracting the 2 neutral entries:
$4\text{ entries}\times 3\times 2 - 2\text{ neutral}\times 3\times 2 = 24 - 12 = 12$.
Both counting schemes agree.
\end{remark}

\begin{resultbox}
Status: CLEAN [L0].  Algebraic matrix argument; no fitting.\\
Independence: No reference to $\alpha$ or $m_e$.\\
Consistency: Agrees with Routes R1, R2, R3.
\end{resultbox}

%──────────────────────────────────────────────────────────────────────────────
\section{Route R5: Compact Mode Counting on $T^3\times S^1_\psi$}
\label{sec:neff12:R5}
%──────────────────────────────────────────────────────────────────────────────

\subsection{Kaluza-Klein Spectrum on the Compactification}
\label{ssec:neff12:R5:KK}

The UBT compactification places the imaginary time on $S^1_\psi$ and spatial
directions on $T^3$.  The Kaluza-Klein (KK) spectrum of $\nabla^\dagger\nabla$
on $T^3\times S^1_\psi$ has eigenvalues:
\begin{equation}
  E_{\mathbf{k},n} = \frac{|\mathbf{k}|^2}{R_t^2} + \frac{n^2}{R_\psi^2},
  \qquad \mathbf{k}\in\mathbb{Z}^3,\; n\in\mathbb{Z}.
  \label{eq:neff12:R5:KK}
\end{equation}

\subsection{First Non-Trivial Winding Modes}
\label{ssec:neff12:R5:first}

\begin{theorem}[$N_{\mathrm{eff}} = 12$ from first KK winding modes]
\label{thm:neff12:R5}
The first non-trivial winding modes on $S^1_\psi$ (those with $|n| = 1$ and
$\mathbf{k} = \mathbf{0}$) include exactly 12 independent charged modes:
\begin{equation}
  N_{\mathrm{eff}}
  = \underbrace{2}_{n=\pm 1}
    \times
    \underbrace{3}_{N_{\mathrm{phases}}}
    \times
    \underbrace{2}_{N_{\mathrm{charge}}}
  = 12.
  \label{eq:neff12:R5}
\end{equation}
\end{theorem}

\begin{proof}
For $|n| = 1$:
\begin{itemize}
  \item $n = +1$ and $n = -1$ are two distinct winding sectors.
  \item Each winding mode exists in three independent Im$\,\mathbb{H}$ phase directions.
  \item Each (winding, phase) combination comes in a charge-conjugate pair.
\end{itemize}
Total: $2 \times 3 \times 2 = 12$ independent modes.

These $|n|=1$ modes dominate the one-loop vacuum polarisation at the
compactification scale because higher winding modes $|n| \geq 2$ are suppressed
by the factor $e^{-|n|/R_\psi}$ in the thermal sum.
\end{proof}

\begin{remark}[Physical interpretation]
The $|n|=1$ winding modes are the analogue of the lightest KK particles in
Kaluza-Klein compactification.  They are the modes that set the effective
coupling at the compactification scale.  Counting exactly these gives $N_{\mathrm{eff}} = 12$.
\end{remark}

\begin{resultbox}
Status: CLEAN [L0].  KK spectrum argument; no fitting.\\
Independence: No reference to $\alpha$ or $m_e$.\\
Consistency: Agrees with Routes R1–R4.
\end{resultbox}

%──────────────────────────────────────────────────────────────────────────────
\section{Convergence and Over-Determination}
\label{sec:neff12:convergence}
%──────────────────────────────────────────────────────────────────────────────

\begin{center}
\begin{tabular}{llll}
\toprule
Route & Starting point & Mechanism & Result \\
\midrule
R1 & $\mathcal{B} = \mathbb{C}\otimes\mathbb{H}$ algebra & Phase $\times$ helicity $\times$ charge & $3\times2\times2 = 12$ \\
R2 & SM gauge group $G_{\mathrm{SM}}$ from $\mathcal{B}$ & Generator count & $8+3+1 = 12$ \\
R3 & 3-qubit internal sectors & Color+isospin+Y $\times$ 2 & $6\times2 = 12$ \\
R4 & $M_2(\mathbb{C})$ matrix structure & Off-diagonal $\times$ phases $\times$ helicity & $2\times3\times2 = 12$ \\
R5 & KK spectrum on $T^3\times S^1_\psi$ & First winding modes & $2\times3\times2 = 12$ \\
\bottomrule
\end{tabular}
\end{center}

\begin{insightbox}
\textbf{Over-determination}: All five routes, starting from different algebraic
perspectives on $\mathcal{B} = \mathbb{C}\otimes\mathbb{H}$, yield $N_{\mathrm{eff}} = 12$.

The common algebraic fact underlying all five routes is:
\begin{equation}
  \dim_\mathbb{R}(\mathrm{Im}\,\mathbb{H}) = 3.
  \label{eq:neff12:common}
\end{equation}
This is the single source of the factor 3 in $N_{\mathrm{eff}} = 12 = 3 \times 4$.

\medskip
\textbf{Non-circularity}: None of the five routes uses $\alpha$ or $m_e$ as input.
The stress test (varying $N_{\mathrm{eff}}$ in \texttt{alpha\_best\_route.tex} \S8)
confirms that only $N_{\mathrm{eff}} = 12$ (from the SM gauge structure in $\mathcal{B}$)
yields $n^* = 137$.
\end{insightbox}

%──────────────────────────────────────────────────────────────────────────────
\section{Connection to the Exponent $3/2$}
\label{sec:neff12:connection}
%──────────────────────────────────────────────────────────────────────────────

The same algebraic factor $N_{\mathrm{phases}} = \dim_\mathbb{R}(\mathrm{Im}\,\mathbb{H}) = 3$
that determines $N_{\mathrm{eff}} = 3 \times 2 \times 2 = 12$ also determines the
exponent $3/2 = N_{\mathrm{phases}}/2$ in $B_{\mathrm{base}} = N_{\mathrm{eff}}^{3/2}$.

This means:
\begin{equation}
  B_{\mathrm{base}} = N_{\mathrm{eff}}^{3/2} = (N_{\mathrm{phases}} \times 4)^{N_{\mathrm{phases}}/2}
  \label{eq:neff12:Bbase_unified}
\end{equation}
is a function of a \emph{single} algebraic input:
$N_{\mathrm{phases}} = \dim_\mathbb{R}(\mathrm{Im}\,\mathbb{H}) = 3$.

This is the deepest structural result: the entire formula
$B_{\mathrm{base}} = N_{\mathrm{eff}}^{3/2}$ reduces to a single fact about the
biquaternion algebra — the dimension of its imaginary quaternion subspace.

\ifdefined\INCLUDEMODE
\else
\end{document}
\fi
